\section{Discussion}\label{sec:disc}

\subsection{What the results permit}\label{sec:disc:uses}

The theorems are statements about one model: a structured nonconvex instance
whose linear costs carry independent random terms of a known scale. Within
that model they permit three things.

First, they give exact global solutions of the sampled instance with expected
bit work bounded in the structural parameter and in explicit numerical
ratios. For problems whose linear costs are themselves uncertain or
estimated, the sampled instance can be the problem of interest, and the
theorems say that its exact global solution is typically cheap under the
stated structure. They do not say that a particular fixed instance is easy.

Second, through \cref{prop:model:regret}, they give a feasible point of the
original problem together with a rational certificate that its original
objective value is within a prescribed $\varepsilon$ of optimal, on every
draw. The expected work is polynomial in $1/\varepsilon$ for fixed structure.
This is an approximation guarantee of the usual kind, obtained by a
different route; we do not claim that its dependence on $\varepsilon$ is
better than that of dedicated approximation schemes.

Third, the proofs isolate certificates of global optimality that are checked
exactly and are sound on every draw, whatever the noise: corrected-corner
bounds on value functions, regions of parametric convex programs on which one
formula is valid, strongly convex patches after sound fixing of integer
labels and active bounds, slab-exclusion certificates obtained from restricted
recourse, competing-label gaps, and optimal-face certificates for convex
integer flows and totally unimodular systems. The probabilistic analysis
concerns only how soon these certificates pass. Such certificates could also
serve as termination tests inside other global solvers, for instance after
branching has localized the optimum; whether they are effective in practice
is not addressed here.

The structural classes cover problems that occur in applications. A small
continuous core with convex, network-flow, or totally unimodular recourse
describes two-stage problems with a few nonconvex first-stage decisions.
Separable convex costs minus a concave term of low rank describe economies
of scale in a few aggregate quantities, with arbitrarily many integer
variables. Order constraints with binary variables describe nonconvex
isotonic fitting with switching, and products of simplices describe
allocation problems. Explicit and implicit graph equalities describe
polynomial or monotone actuator models. In each case the theorems apply
exactly as stated, including their numerical ratios and oracle assumptions;
we have not evaluated any of these classes computationally.

The results also indicate where randomness helps. Noise on the coefficients
that carry the nonconvexity is what the counts use. Aligned and Gaussian-like
noise give dimension-free counts in the low-rank route, while uniform ambient
noise does not (\cref{thm:lim:ambient}). Noise on a small core suffices to
locate the core, but residual point output then needs residual strong
convexity or a certificate for all tied residual solutions
(\cref{sec:lim:points}). Noise on every coordinate supplies point growth and
active-bound margins in the sparse and ambient recourse routes.

\subsection{Limitations}\label{sec:disc:limits}

The bounds depend on numerical ratios such as $L/\sigma$ or
$\nu\diam(X)/\sigma$ through their magnitudes. Small noise makes them large,
and a ratio that is exponential in the input length gives an exponential
bound. The sparse bounds are polynomial only for fixed width, and their
dimension factor is necessary for the retention rule used
(\cref{thm:lim:width}).

The finite law is constructed from the base instance after a fallback cost
factor and a stopping depth are known. The analysis does not provide one
universal law for all instances of a given input length, and it does not
cover coarse laws chosen independently of the instance. Ties at atoms are
handled by an exact fallback whose worst-case cost is exponential; only its
expected contribution is polynomial.

Several premises are supplied rather than discovered: curvature bounds that
are sharper than coefficient bounds, a core, a tree decomposition, a
quadratic convexifier in the supplied-factor variants, residual convexity or
a residual strong-convexity modulus, and exact or certified recourse oracles
that remain valid under every coordinate restriction. Where a premise is
promised rather than verified, the theorem says whether correctness or only
the work bound depends on it. Verifying a general polynomial inequality is
itself hard, and the theorems do not hide such verification in the input
length.

The outputs are exact but not always explicit. Implicit descriptors have
polynomial length, and evaluation to any precision is efficient; exact
expanded algebraic coordinates, exact comparison of the optimal value with
a threshold, and the exact active set are not provided. Fallback outputs can
be exponentially long on rare draws, and the global proof records have
polynomial size only in expectation.

The constants are large. The fallback factor is exponential in a polynomial
of the input length, the component solver has a fixed but astronomically
large base, and the sampling precision, although polynomial, is large. The
results are therefore mathematical guarantees, not predictions of practical
running time. No implementation or experiment accompanies them, and no
practical speedup is claimed.

Finally, every theorem solves the sampled objective. The original objective
receives only the certified approximation of \cref{prop:model:regret}; the
results do not solve unperturbed instances exactly, which is consistent with
their worst-case hardness.

\subsection{Open questions}\label{sec:disc:open}

\emph{Width.} Is there an exact algorithm for sparse polynomial optimization
under independent linear noise whose expected work is fixed-parameter in the
width? \Cref{thm:lim:width} and \cref{prop:lim:local} show that this needs
certified conditional values for the variables outside a bag, with error
controlled by the bag alone; grid min-marginals do not provide them. The
recourse route provides such values for a small core, which suggests looking
for decompositions in which every bag has a tractable conditional problem.

\emph{Uniform ambient noise.} Is there a fixed-parameter bound in the negative
inertia under independent uniform noise on all coefficients?
\Cref{thm:lim:ambient} shows that it would need a different argument than
counting near-optimal auxiliary grid nodes.

\emph{Constraints.} Which classes of sparse linear constraints, beyond
products of simplices, order constraints, and graph equalities, admit exact
smoothed algorithms under objective noise? \Cref{thm:lim:constraints} shows
that some structure beyond sparsity is necessary, and
\cref{ex:lim:coupled} shows that coordinatewise counting must be replaced by
counts adapted to the faces of the feasible set.

\emph{Residual points under core-only noise.} With merely convex residual
problems, \cref{thm:lim:posslp} shows that efficient evaluation of residual
points would give a Las Vegas algorithm for $\PosSLP$, even on instances
whose optimal value is easy to approximate. Which additional structure,
short of a uniform strong-convexity modulus, suffices for residual point
output? Certificates that select an optimal residual face, as the flow
certificate of \cref{sec:int} does for integer recourse, are one candidate.

\emph{Universal laws and exact formats.} Can the instance-dependent finite
law be replaced by one law for all instances of a given input length, with
the same expected bounds? Which structured classes admit exact algebraic
outputs of polynomial size on every draw, rather than implicit descriptors
and expected-size fallback outputs?

\emph{Computation.} The certificates are exact and checkable, but their
practical value depends on certified convex solvers, compact proof records,
and much smaller constants than those proved here. Implementations and
computational studies are needed before any practical conclusion can be
drawn.
